% !TeX root = ../main.tex
\chapter{Relazioni, equivalenze e quozienti}
\fonti{Lezioni 2, 3 e 4}

\section{Relazioni tra insiemi}
% A–L L2 pp. 3–4; M–Z L2 p. 7 e L3 pp. 1–3.
\begin{definizione}
Una \textbf{relazione} da $X$ a $Y$ è un sottoinsieme $\rho\subseteq X\times Y$. La scrittura $x\rho y$ significa $(x,y)\in\rho$. Se i due insiemi coincidono, si parla di relazione binaria \textbf{su} $X$.
\end{definizione}

\begin{esempio}[Città visitate]
Siano $X=\{P,V,R,M,T\}$ cinque città e $Y=\{\text{Anna},\text{Luca},\text{Lia}\}$. La relazione «la città è stata visitata dalla persona» può essere
\[
\begin{split}
\rho=\{&(P,\text{Anna}),(V,\text{Anna}),(R,\text{Anna}),\\
&(M,\text{Luca}),(T,\text{Luca}),(M,\text{Lia}),(R,\text{Lia})\}.
\end{split}
\]
La rappresentazione per tabella riporta le stesse coppie:
\[
\begin{array}{c|ccc}
&\text{Anna}&\text{Luca}&\text{Lia}\\\hline
P&\times&&\\V&\times&&\\R&\times&&\times\\M&&\times&\times\\T&&\times&
\end{array}
\]
Una città può essere in relazione con più persone e una persona con più città; una relazione non richiede alcuna unicità.
\end{esempio}

% M–Z L2 p. 7 e L3 pp. 1–2: esempio distinto, persone per città e righe/colonne vuote.
\begin{esempio}[Persone e città: nessuna condizione di esistenza]
Consideriamo invece una relazione dall'insieme delle persone $U$ all'insieme delle città $V$, dove
\[
\begin{aligned}
U&=\{\text{Lidia},\text{Marco},\text{Pietro},\text{Francesca}\},\\
V&=\{\text{Roma},\text{Parigi},\text{Madrid},\text{Lisbona},\text{Vienna}\}.
\end{aligned}
\] Lidia ha visitato Roma e Madrid; Marco Vienna e Parigi; Francesca Roma, Parigi e Madrid; Pietro nessuna delle cinque. Le sette coppie della relazione $\eta\subseteq U\times V$ sono quindi
\[
\begin{split}
\eta=\{&(\text{Lidia},\text{Roma}),(\text{Lidia},\text{Madrid}),\\
&(\text{Marco},\text{Vienna}),(\text{Marco},\text{Parigi}),\\
&(\text{Francesca},\text{Roma}),(\text{Francesca},\text{Parigi}),\\
&(\text{Francesca},\text{Madrid})\}.
\end{split}
\]
Qui la prima coordinata è una persona e la seconda una città, diversamente dall'esempio precedente. Un elemento del primo insieme può non essere in relazione con alcun elemento del secondo (Pietro), e un elemento del secondo può non comparire in alcuna coppia (Lisbona). Una relazione non impone né esistenza né unicità di un elemento associato.
\end{esempio}

\begin{esercizio}[Grafici di relazioni]
Rappresentare $\rho=\{(x,y)\in\N\times\Z\mid x+y=1\}$, $\sigma=\{(x,y)\in\R^2\mid x^2+y^2=1\}$ e $\tau=\{(a,b)\in\N^*\times\N\mid b/a\in\N\}$.
\end{esercizio}
\begin{svolgimento}
Per $\rho$ si ottengono i punti $(n,1-n)$, $n\in\N$, a partire da $(0,1)$: è una parte discreta della retta $y=1-x$, non tutta la retta. Il grafico di $\sigma$ è la circonferenza di centro l'origine e raggio $1$. Per $\tau$, fissata un'ascissa intera $a>0$, le ordinate sono $b=ka$, $k\in\N$: tutti i multipli non negativi di $a$, compreso $0$. La scrittura $a\mid b$, letta «$a$ divide $b$», significa che $b=ak$ per un intero $k$; per i valori qui considerati tale intero è non negativo.
\end{svolgimento}

\begin{schema}
\begin{tikzpicture}[scale=.65]
  \begin{scope}
    \draw[asse] (-.5,0)--(4,0) node[right] {$x$};\draw[asse] (0,-2.6)--(0,1.8) node[above] {$y$};
    \draw[densely dashed,black!45] (-.3,1.3)--(3.4,-2.4);
    \foreach \x in {0,...,3}{\fill[coloreschema] (\x,{1-\x}) circle (2pt);}
    \node at (1.6,-3.2) {$x+y=1$, $x\in\N$};
  \end{scope}
  \begin{scope}[shift={(7,0)}]
    \draw[asse] (-1.7,0)--(1.8,0) node[right] {$x$};\draw[asse] (0,-1.6)--(0,1.8) node[above] {$y$};
    \draw[coloreschema,thick] (0,0) circle (1);\node at (.9,.3) {$1$};
    \node at (0,-3.2) {$x^2+y^2=1$};
  \end{scope}
  \begin{scope}[shift={(11,0)}]
    \draw[asse] (-.3,0)--(4.5,0) node[right] {$a$};\draw[asse] (0,-.3)--(0,4.6) node[above] {$b$};
    \foreach \a in {1,...,4}{\foreach \b in {0,...,4}{\pgfmathtruncatemacro{\rem}{mod(\b,\a)}\ifnum\rem=0\fill[coloreschema] (\a,\b) circle (2pt);\fi}}
    \node at (2,-3.2) {$a\mid b$, porzione del grafico};
  \end{scope}
\end{tikzpicture}
\end{schema}

\section{Equivalenze e ordini}
\begin{definizione}
Una relazione $\rho$ su $X$ è \textbf{riflessiva} se $x\rho x$ per ogni $x\in X$; \textbf{simmetrica} se $x\rho y\Rightarrow y\rho x$ per ogni $x,y\in X$; \textbf{transitiva} se $x\rho y$ e $y\rho z$ implicano $x\rho z$ per ogni $x,y,z\in X$. Una \textbf{relazione di equivalenza} soddisfa tutte e tre le proprietà.
\end{definizione}

\begin{esempio}[Uguaglianza]
La diagonale $\Delta_X=\{(x,x)\mid x\in X\}$ definisce l'uguaglianza: $x=x$; da $x=y$ segue $y=x$; da $x=y$ e $y=z$ segue $x=z$. Sono precisamente le tre proprietà richieste.
\end{esempio}

\begin{esercizio}[Tre equivalenze]
Verificare che sono equivalenze la relazione $(a,b)\sim(c,d)\iff ad=bc$ su $F=\Z\times\Z^*$, la similitudine tra triangoli non degeneri del piano e il parallelismo tra rette, ammettendo rette coincidenti.
\end{esercizio}
\begin{svolgimento}
\textbf{Coppie intere.} La riflessività è $ab=ba$. La simmetria segue da $ad=bc\Rightarrow cb=da$. Per la transitività, se $(a,b)\sim(c,d)$ e $(c,d)\sim(e,f)$, allora $ad=bc$ e $cf=de$. Moltiplicando si ottiene $adf=bcf=bde$; poiché $d\ne0$, si cancella $d$ e si ricava $af=be$, cioè $(a,b)\sim(e,f)$.

\textbf{Triangoli.} Per il criterio degli angoli, due triangoli sono simili se i rispettivi angoli si possono mettere in corrispondenza con ampiezze uguali. Ogni triangolo ha gli stessi angoli di se stesso; una corrispondenza di angoli uguali vale anche in senso contrario; se gli angoli del primo coincidono con quelli del secondo e quelli del secondo con quelli del terzo, coincidono anche quelli del primo e del terzo. Si hanno riflessività, simmetria e transitività.

\textbf{Rette.} Essere parallele, ammettendo la coincidenza, significa avere la stessa direzione. L'uguaglianza delle direzioni è riflessiva, simmetrica e transitiva. Escludendo la coincidenza si perderebbe la riflessività.
\end{svolgimento}

\begin{esercizio}[Riconoscere le proprietà]
Stabilire quali delle seguenti relazioni sono equivalenze: $a\rho b\iff b=a^2$ su $\R$; la perpendicolarità tra rette; $a\sigma b\iff a\le b$ su $\N$; $a\tau b\iff |a|=|b|$ su $\R$; $x\sim_f y\iff f(x)=f(y)$, dove a ogni elemento $x$ è associato un unico valore $f(x)$.
\end{esercizio}
\begin{svolgimento}
\textbf{Quadrato.} Non è riflessiva ($2\not\rho2$), né simmetrica ($2\rho4$, ma $4\not\rho2$), né transitiva ($2\rho4$ e $4\rho16$, ma $2\not\rho16$).

\textbf{Perpendicolarità.} È simmetrica, ma una retta non è perpendicolare a se stessa. Non è transitiva: da $r\perp s$ e $s\perp r$ non segue $r\perp r$.

\textbf{Ordine usuale.} È riflessivo e transitivo, ma $0\le1$ e $1\nleq0$, quindi non simmetrico.

\textbf{Valore assoluto e uguaglianza dei valori.} Sono entrambe equivalenze: si applicano le tre proprietà dell'uguaglianza, rispettivamente ai numeri $|a|$ e ai valori $f(x)$. In particolare $f(x)=f(y)=f(z)$ implica $f(x)=f(z)$, senza richiedere che $x=y=z$.
\end{svolgimento}

\begin{definizione}
Una relazione è \textbf{antisimmetrica} se $x\rho y$ e $y\rho x$ implicano $x=y$. Un \textbf{ordine parziale} è una relazione riflessiva, antisimmetrica e transitiva. È un \textbf{ordine totale} se ogni coppia è confrontabile: per ogni $x,y$ vale $x\rho y$ oppure $y\rho x$.
\end{definizione}

\begin{esercizio}[Esempi di ordini]
Verificare che l'inclusione è un ordine parziale su $\mathcal P(X)$, la divisibilità è un ordine parziale su $\N^*$ e $\le$ è un ordine totale su $\R$. Che cosa cambia per la divisibilità su $\Z^*$?
\end{esercizio}
\begin{svolgimento}
\textbf{Inclusione.} $A\subseteq A$; $A\subseteq B\subseteq A$ implica $A=B$; $A\subseteq B\subseteq C$ implica $A\subseteq C$. Le tre proprietà seguono direttamente dall'appartenenza.

\textbf{Divisibilità.} Scriviamo $a\mid b$ se $b=ak$ per qualche intero $k$. Si ha $a=a\cdot1$. Se $b=ak$ e $c=bh$, allora $c=a(kh)$, quindi vale la transitività. Per $a,b>0$, da $a\mid b$ segue $a\le b$; se anche $b\mid a$, si ha $a=b$. Non è totale: $2$ e $3$ non sono confrontabili. Su $\Z^*$ viene meno l'antisimmetria: $1\mid-1$ e $-1\mid1$, ma $1\ne-1$.

\textbf{Ordine reale.} Le proprietà usuali di $\le$ danno riflessività, antisimmetria, transitività e confrontabilità di ogni coppia, quindi un ordine totale.
\end{svolgimento}

\begin{esercizio}[Confrontabilità]
Per $X\ne\varnothing$, l'inclusione su $\mathcal P(X)$ è sempre totale? Su $S=\{1,2,-2,6,-6,18\}$, con la divisibilità intera, ogni coppia è confrontabile? Vale l'antisimmetria?
\end{esercizio}
\begin{svolgimento}
Se $X$ contiene $x\ne y$, i singoletti $\{x\}$ e $\{y\}$ non sono confrontabili. Se $X$ ha un solo elemento, invece, $\mathcal P(X)=\{\varnothing,X\}$ è una catena.

In $S$ i valori assoluti sono ordinati dalla divisibilità come $1\mid2\mid6\mid18$. Il segno non modifica la divisibilità tra interi, perciò ogni coppia è confrontabile. Ma $2\mid-2$ e $-2\mid2$ con $2\ne-2$: confrontabilità e antisimmetria sono condizioni diverse, e qui non si ha un ordine totale.
\end{svolgimento}

\section{Partizioni, classi e quozienti}
% A–L L3 pp. 1, 3–7; M–Z L3 pp. 5–9. Conservata la precedenza A–L della partizione.
\begin{definizione}
Una \textbf{partizione} di $X$ è un insieme $\mathcal A$ di sottoinsiemi non vuoti di $X$, detti \textbf{blocchi}, la cui unione è $X$ e tali che due blocchi distinti siano disgiunti. Equivalentemente,
\[
\bigcup_{A\in\mathcal A}A=X,\qquad
A\cap B\ne\varnothing\ \Longleftrightarrow\ A=B
\quad(A,B\in\mathcal A).
\]
Nella seconda formulazione, prendendo $A=B$, si ricava anche che ogni blocco è non vuoto. Ogni elemento di $X$ appartiene a un unico blocco. Per $X=\varnothing$ la partizione è la famiglia vuota.
\end{definizione}

\begin{esempio}[Scomposizione di un'unione]
I tre pezzi $A\setminus B$, $A\cap B$, $B\setminus A$ ricoprono $A\cup B$ e sono a due a due disgiunti; in particolare
\[
A=(A\setminus B)\cup(A\cap B),\qquad
B=(B\setminus A)\cup(A\cap B).
\]
Eliminando i pezzi vuoti si ottiene una partizione di $A\cup B$.

Per tre insiemi, richiamiamo le regioni $A_0=A\setminus(B\cup C)$, $B_0=B\setminus(A\cup C)$, $C_0=C\setminus(A\cup B)$ e $T=A\cap B\cap C$. Le altre tre sono
\[
R_{AB}=(A\cap B)\setminus C,\quad
R_{AC}=(A\cap C)\setminus B,\quad
R_{BC}=(B\cap C)\setminus A.
\]
Le sette regioni sono a due a due disgiunte e ricoprono $A\cup B\cup C$: ogni punto appartiene a esattamente uno, due oppure tutti e tre gli insiemi, e il suo elenco di appartenenze individua la regione. In particolare,
\[
\begin{aligned}
A&=A_0\cup R_{AB}\cup R_{AC}\cup T,\\
B&=B_0\cup R_{AB}\cup R_{BC}\cup T,\\
C&=C_0\cup R_{AC}\cup R_{BC}\cup T.
\end{aligned}
\]
Eliminando le regioni vuote si ottiene una partizione dell'unione. Per esempio $A\triangle B=A_0\cup B_0\cup R_{AC}\cup R_{BC}$; facendo poi la differenza simmetrica con $C$ rimangono $A_0,B_0,C_0,T$, come nella dimostrazione dell'associatività.
\end{esempio}

\begin{definizione}
Per un'equivalenza $\sim$ su $X$, la \textbf{classe di equivalenza} di $x\in X$ è $[x]=\{y\in X\mid y\sim x\}$. Ogni suo elemento è un \textbf{rappresentante} della classe. L'\textbf{insieme quoziente} è $X/{\sim}=\{[x]\mid x\in X\}\subseteq\mathcal P(X)$.
\end{definizione}

\begin{proposizione}\label{prop:classi}
Per $x,y\in X$, $x\sim y$ se e solo se $[x]=[y]$.
\end{proposizione}
\begin{proof}
\passomath{x\sim y\Rightarrow[x]\subseteq[y]} Se $z\in[x]$, allora $z\sim x$ e $x\sim y$; per transitività $z\sim y$, quindi $z\in[y]$.

\passomath{x\sim y\Rightarrow[y]\subseteq[x]} Per simmetria $y\sim x$. Se $z\in[y]$, da $z\sim y\sim x$ segue $z\in[x]$. Le due inclusioni danno l'uguaglianza.

\passomath{[x]=[y]\Rightarrow x\sim y} Per riflessività $x\in[x]$; dall'uguaglianza segue $x\in[y]$, che significa precisamente $x\sim y$.
\end{proof}

\begin{teorema}\label{teo:partizioni}
\begin{enumerate}[label=\arabic*.]
\item Le classi di un'equivalenza formano una partizione di $X$.
\item Viceversa, se $\mathcal A$ è una partizione, la relazione $x\sim_{\mathcal A}y$ definita dall'appartenenza a uno stesso blocco è un'equivalenza, le cui classi sono i blocchi di $\mathcal A$.
\end{enumerate}
\end{teorema}
\Needspace{4\baselineskip}
\begin{proof}
\begin{enumerate}[label=\arabic*.]
\item \textbf{Dall'equivalenza alla partizione.} Ogni classe contiene il proprio rappresentante, dunque non è vuota. Ogni $x\in X$ appartiene a $[x]$, perciò l'unione delle classi contiene $X$; l'altra inclusione vale perché ogni classe è contenuta in $X$. Se $z\in[x]\cap[y]$, si ha $z\sim x$ e $z\sim y$. Per simmetria e transitività segue $x\sim z\sim y$, quindi $[x]=[y]$. Due classi distinte sono pertanto disgiunte.
\item \textbf{Dalla partizione all'equivalenza.} Ogni $x$ appartiene a un blocco, dunque $x\sim_{\mathcal A}x$. Se $x,y$ stanno nello stesso blocco, vi stanno anche $y,x$, quindi la relazione è simmetrica. Se $x,y\in A$ e $y,z\in B$, i blocchi si incontrano in $y$ e perciò $A=B$; allora $x,z$ stanno nello stesso blocco, ottenendo la transitività.

\textbf{Identificazione delle classi.} Sia $x\in A\in\mathcal A$. Ogni $y\in A$ è equivalente a $x$, quindi $A\subseteq[x]$. Viceversa, se $y\sim_{\mathcal A}x$, un blocco $B$ contiene entrambi; poiché $x\in A\cap B$, si ha $A=B$ e $y\in A$. Dunque $[x]=A$.
\end{enumerate}
\end{proof}

\begin{schema}
\begin{tikzpicture}
  \draw[rounded corners] (0,0) rectangle (6,2.4);
  \draw (2,0)--(2,2.4);\draw (4,0)--(4,2.4);
  \node at (1,1.9) {$[x]$};\node at (3,1.9) {$[y]$};\node at (5,1.9) {$[z]$};
  \foreach \x/\y in {.5/.7,1.4/1,2.6/.6,3.4/1.1,4.6/.7,5.3/1.1}{\fill[coloreschema] (\x,\y) circle (1.5pt);}
  \node[right] at (6.2,1.2) {$X$};
\end{tikzpicture}
\end{schema}

\begin{esercizio}[Le costruzioni inverse]
Verificare che le due costruzioni del teorema stabiliscono una corrispondenza uno a uno tra equivalenze su $X$ e partizioni di $X$.
\end{esercizio}
\begin{svolgimento}
Partendo da una partizione, il teorema mostra che le classi della relazione costruita sono esattamente i blocchi di partenza: si recupera la partizione.

Partendo da $\sim$, due elementi $x,y$ appartengono a una stessa classe se e solo se $[x]=[y]$, cioè, per la proposizione~\ref{prop:classi}, se e solo se $x\sim y$. Si recupera anche la relazione iniziale. Ciascuna costruzione annulla dunque l'altra.
\end{svolgimento}

\begin{esercizio}[Descrivere i quozienti]
Descrivere le classi e i quozienti degli esempi di equivalenza precedenti, compresa la relazione $x^2=y^2$ su $\R$.
\end{esercizio}
\begin{svolgimento}
Per l'uguaglianza, le classi sono i singoletti $\{x\}$. Per $|x|=|y|$, sono $\{0\}$ e $\{t,-t\}$ con $t>0$; lo stesso vale per $x^2=y^2$, perché $x^2-y^2=(x-y)(x+y)$. Per esempio $[5]=[-5]=\{5,-5\}$.

Per il parallelismo, una classe comprende tutte le rette di una stessa direzione. Per la similitudine, comprende tutti i triangoli di una stessa forma: la terna degli angoli, senza fissarne l'ordine, determina la classe. Per l'uguaglianza $f(x)=f(y)$, una classe raccoglie tutti gli elementi ai quali è associato lo stesso valore. Il quoziente delle coppie intere per prodotti incrociati è descritto di seguito.
\end{svolgimento}

\section{I razionali e le operazioni sul quoziente}
% A–L L3 p. 8 e L4 pp. 1–3; M–Z L4 pp. 1–4.
\begin{definizione}
Posto $F=\Z\times\Z^*$ e $(a,b)\sim(c,d)\iff ad=bc$, si definisce $\Q=F/{\sim}$. La notazione $a/b$ indica la classe $[(a,b)]$, non la singola coppia. Per esempio
\[
\frac{-2}{4}=\frac{-1}{2}=\frac{1}{-2},\qquad
\frac12=\frac24=\frac{-2}{-4}=\frac36.
\]
\end{definizione}

Identifichiamo un intero $n$ con la classe $n/1$: interi distinti restano distinti, perché $n/1=m/1$ equivale a $n=m$ per il criterio dei prodotti incrociati.

Una regola sui rappresentanti dà un'operazione sulle classi solo se il risultato è indipendente dai rappresentanti scelti.

\begin{esempio}[Una regola incompatibile]
La formula $a/b\star c/d=(a+c)/(bd)$ non definisce un'operazione su $\Q$: infatti $1/2=2/4$, ma
\[
\frac12\star\frac23=\frac36=\frac12,\qquad
\frac24\star\frac23=\frac4{12}=\frac13.
\]
La stessa coppia di razionali produrrebbe due risultati diversi.
\end{esempio}

% Mantiene insieme enunciato e sviluppo algebrico sul raccordo finale.
\Needspace{30\baselineskip}
\begin{proposizione}[Somma e prodotto]
Le seguenti formule definiscono operazioni su $\Q$:
\begin{enumerate}[label=\arabic*.]
\item $\displaystyle\frac ab+\frac cd=\frac{ad+bc}{bd}$.
\item $\displaystyle\frac ab\cdot\frac cd=\frac{ac}{bd}$.
\end{enumerate}
\end{proposizione}
\begin{proof}
Siano $ab'=a'b$ e $cd'=c'd$, con $b,b',d,d'\ne0$. I denominatori $bd$ e $b'd'$ sono non nulli, dunque le coppie risultanti sono ammesse.

\begin{enumerate}[label=\arabic*.]
\item \textbf{Somma.} Per l'equivalenza delle coppie occorre $(ad+bc)b'd'=(a'd'+b'c')bd$. Le ipotesi danno
\[
\begin{aligned}
(ad+bc)b'd'&=(ab')dd'+bb'(cd')\\
&=(a'b)dd'+bb'(c'd)\\
&=(a'd'+b'c')bd.
\end{aligned}
\]
\item \textbf{Prodotto.} Analogamente $(ac)b'd'=(ab')(cd')=(a'b)(c'd)=(a'c')bd$. Entrambi i risultati appartengono quindi alla stessa classe dopo il cambio di rappresentanti.
\end{enumerate}
\end{proof}

\begin{proposizione}[Criterio di compatibilità]
Siano $\sim$ un'equivalenza su $X$ e $\star$ un'operazione su $X$, cioè una regola che a ogni coppia di elementi di $X$ associa un elemento di $X$. La formula $[x]\mathbin{\overline\star}[y]=[x\star y]$ definisce un'operazione su $X/{\sim}$ se e solo se
\[
x\sim x',\ y\sim y'\quad\Longrightarrow\quad x\star y\sim x'\star y'
\]
per ogni $x,x',y,y'\in X$.
\end{proposizione}
\begin{proof}
\passomath{\Rightarrow} Se $x\sim x'$ e $y\sim y'$, le classi in ingresso coincidono. Un'operazione ben definita deve dare la stessa classe in uscita, ossia $[x\star y]=[x'\star y']$; per la proposizione~\ref{prop:classi} i rappresentanti sono equivalenti.

\passomath{\Leftarrow} Due scelte dei rappresentanti di $[x]$ e $[y]$ differiscono per le equivalenze $x\sim x'$ e $y\sim y'$. La condizione assicura che i risultati abbiano la stessa classe. La regola è quindi indipendente da entrambe le scelte.
\end{proof}
